\begin{statement}{E.43}{Proved}{The double-cover refinement and every retained boundary tower coexist}
Use $G_L,C,T_L,S_L$ from E.38 and the actual subgroup
$H\lhd F\lhd B$ from E.42.  Put $S=G_L/C$.
The literal quotient map from $B$ has the factorization
\[
 B\longrightarrow P_{20}=B/\bigl(\gamma_3(H)H^{20}\bigr)
                  \longrightarrow S,
 \tag{E.43.1}
\]
and its exact $H$ image is
\[
                 Q_8\times C_{10}^2.
 \tag{E.43.2}
\]
The group $S$ is the maximal solvable quotient of $G_L$.

Fix any of the $432$ multiplier-one points of E.24,
any full two-adic class-two fibre in E.41, and any
full three-adic class-two fibre in E.37.  There is one
profinite commutator word retaining all these images,
inducing the identity on every full relative pro-$p$
boundary quotient $\mathscr P_p$ for $p\ge5$, and
satisfying the complete SL$_2(\F_5)$ Ferrari test
of E.38.  All the resulting finite common source
maps are surjective.  The four-strand map onto
$T_{23}\times_{S_L}T_L$ is surjective and satisfies
the full pentagon and twist identities.

More precisely, for every finite old common source
$G_{\rm old}$ in E.42 which retains $P_{20}$, the
actual new source and its corrected image are both
\[
                    G_{\rm old}\times_S G_L.
 \tag{E.43.3}
\]
These old sources are cofinal in the specified
boundary family.  This is a simultaneous extension
of that family, with its actual comparison maps.
\end{statement}
\begin{proof}
We first establish the finite comparison which
resolves the additional common five-primary quotient.
Write $\overline F=\rho_L(F)/C\subset S$ and
$\overline E=\rho_L(F')/C$.
E.38 gives
\[
 |S|=43200,\qquad |\overline F|=7200,
 \qquad \overline E\simeq Q_8,
 \qquad |\overline F^{\rm ab}|=900.
 \tag{E.43.4}
\]
The index six in the second equality is the actual
outer permutation map $B/F\simeq S_3$.

In the faithful $1980$-point $A_5$ source model,
the literal words
$x=(r^{-1}s)^2$ and $y=(sr^{-1})^2$
both have order $15$.
The new projection $G_L\to G_A$ has elementary
abelian two-kernel.  Thus their orders in $G_L$
divide $30$.  Their signed modular images have
order six, so the two actual projections show
that both orders are exactly $30$.
Consequently the images of $x,y$ give an onto
map $(\mathbb Z/30)^2\to\overline F^{\rm ab}$.
Both groups have order $900$, so this is an
isomorphism with these specified generators.
In particular its inverse composed with
$F\to\overline F$ is exactly the ordinary
pair of free exponent sums modulo $30$.

The signed modular map is also specified on
all of $F$.  Its generated image in
$\operatorname{SL}_2(\mathbb Z/12)$ has order
$96$ and contains no pair $g,-g$.
It is therefore the unique lift in that image
of the projective $F$ image used in E.38.
Reduction modulo three gives a homomorphism
$q_3:\overline F\to\operatorname{SL}_2(\F_3)$.
It is injective on $\overline E$ by (E.38.12).
It follows that the map
\[
 \Theta:\overline F\longrightarrow
       \operatorname{SL}_2(\F_3)\times(\mathbb Z/30)^2,
 \quad
 w\longmapsto
 \bigl(q_3(w),\operatorname{ab}_x(w),
                         \operatorname{ab}_y(w)\bigr)
 \tag{E.43.5}
\]
is injective.  Indeed an element in its kernel
first lies in $\overline E$ by abelianization,
and then is identity by the injectivity of $q_3$.

Let $\kappa_3:\operatorname{SL}_2(\F_3)\to C_3$
be its abelianization, chosen to send both actual
matrices $q_3(x),q_3(y)$ to $1$.
The matrix enumeration verifies this map on
all $24$ elements.  Formula \textup{(E.43.5)}
has image precisely
\[
 \left\{(g,a,b):\kappa_3(g)=a+b\pmod3\right\}.
 \tag{E.43.6}
\]
The two generator images lie in this fibre
product, and its order is $24\cdot900/3=7200$,
equal to the proved order of the injected group.
This proves equality of the groups with their
specified maps from $F$.
The program also enumerates all $7200$ values
directly from these two generator images.

The defining map $F\to F/H=(\mathbb Z/3)^2$
is the reduction of the last two coordinates
in \textup{(E.43.5)}.
It follows that the kernel of $B\to S$ lies
in $H$, and that the actual $H$ image is
\[
 \left\{(g,3a,3b):g\in Q_8,
                         \ a,b\in\mathbb Z/10\right\}.
 \tag{E.43.7}
\]
Its multiplication is the direct-product
multiplication, proving \textup{(E.43.2)}
and the exact index $43200/800=54$.
This also identifies $S/(\rho_L(H)/C)$ with
the actual quotient $B/H$.

For an independent centrality comparison,
the actual matrices $q_{12}(x^3)$ and
$q_{12}(y^3)$ commute with the entire
$96$-element modular $F$ image.
Their commutators with $\overline F$ lie in
$\overline E$, where $q_{12}$ is injective.
Thus $x^3,y^3$ are central in $\overline F$.
They generate $C_{10}^2$ by their independent
abelianization coordinates, intersecting
$\overline E$ trivially.  This recovers
\textup{(E.43.7)} using the original source
generators directly.

The certificate retains all ten literal
Schreier generators of $H$, from the section
$x^a y^b$, $0\le a,b<3$, and their values
in \textup{(E.43.7)}.
Their images generate all $800$ elements.
The two actual Sylow projections are
\[
 h\longmapsto h^5\in Q_8\times C_2^2,
 \qquad
 h\longmapsto h^{16}\in C_5^2.
 \tag{E.43.8}
\]
The images have orders $32$ and $25$, commute,
and multiply back to $h$ for all $800$ values.
In the coordinates of \textup{(E.43.7)}
they are the reductions of $(a,b)$ modulo
two and five, retaining $g$ in the first
factor.  The five-primary map is explicitly
\[
 h\longmapsto
 \bigl(\operatorname{ab}_x(h)/3,
       \operatorname{ab}_y(h)/3\bigr)\pmod5;
 \tag{E.43.9}
\]
division is legitimate on $H$, and is
equivalently multiplication by two of
the ordinary exponent sums modulo five.
Every third lower-central commutator and
every twentieth power vanish in this
finite $H$ image; its derived group has
order two and its exponent is exactly
$20$.  These identities prove
\textup{(E.43.1)} for the literal images
of $s,r$.  The subgroups in its kernel
are characteristic in $H$ and hence
normal in $B$, so this is a homomorphism
of the displayed actual quotient groups.

The group in \textup{(E.43.7)} is solvable,
and $B/H$ is the extension of the abelian
group $F/H$ by $S_3$.  Hence $S$ is
solvable.  Conversely every solvable
quotient of $G_L$ kills the perfect
normal subgroup $C$.  Thus $S$ is the
maximal solvable quotient, with the
specified map from $B$.

\paragraph{The old system contains this quotient.}
The actual Hall-coordinate maps give
\[
 P_{20}\xrightarrow{\sim}
                  P_4\times_{B/H}P_5.
 \tag{E.43.10}
\]
On the geometric kernel, the abelian
coordinates are modulo $20$ and the
central Hall coordinates modulo $10$.
Indeed the power and commutator formulas
give the central modulus
$\gcd(20,\binom{20}{2})=10$;
conversely these Hall moduli satisfy
every twentieth-power relation.
The maps on these coordinates are the
Chinese-remainder isomorphisms to
moduli $(4,5)$ and $(2,5)$, respectively.
They respect collection because they
are the natural quotient homomorphisms
of $H$.  The two geometric factors
have coprime orders, so the image of
$H$ is their full product.  Together
with the common quotient $B/H$, this
proves \textup{(E.43.10)}.
E.41 retains $P_4$, and E.42 retains
$P_5$ as a quotient of its full
relative pro-five factor.
Consequently adjoining this fixed
$P_{20}$ to any old finite source gives
a cofinal subsystem of the family in
E.42, with an actual map to $S$ by
\textup{(E.43.1)}.

\paragraph{A correction preserving every solvable boundary image.}
Let $f_0\in\widehat H\cap\widehat F'$
be the actual word supplied by E.42
for the prescribed boundary parameters.
It retains the chosen $G_6$ point, so
its $q_{12}$ value is still the indexed
matrix in (E.38.13).
Choose an actual Galois element
$\sigma$ with $\chi(\sigma)=1$ and
the corresponding Kummer value, as
in E.38.  The exact quotient comparison
then gives
\[
             c=\rho_L(f_0)^{-1}\rho_L(f_\sigma)\in C.
 \tag{E.43.11}
\]

For a profinite group $J$, let $J^{[r]}$
be its closed derived series, beginning
with $J^{[0]}=J$.
Since $\rho_L(\widehat N^{[1]})=C$
and $C$ is perfect,
\[
             \rho_L(\widehat N^{[r]})=C
                              \qquad(r\ge1).
 \tag{E.43.12}
\]
The preimages of the fixed $c$ inside
these compact groups are nonempty,
nested closed sets.  Their intersection
therefore gives an element $k$ with
\[
 k\in\bigcap_{r\ge1}\widehat N^{[r]},
                 \qquad \rho_L(k)=c.
 \tag{E.43.13}
\]
Here $N'\subset F'\subset H$.
Induction on the derived series gives
$\widehat N^{[r+1]}\subset\widehat H^{[r]}$.
Thus $k$ belongs to the full prosolvable
residual $\bigcap_r\widehat H^{[r]}$.
Every homomorphism from $\widehat H$
to a finite solvable group kills this
intersection, by the finite derived
length of its image.
This applies to the $H$ image of
$\mathscr D$, to every finite two- or
three-adic class-two boundary image,
and to every finite $p$-group quotient
of $\widehat H$ for every $p\ge5$.
Consequently $k$ is identity in every
component retained in E.42, simultaneously.
More explicitly, this proves the exact
continuous surjection
\[
 R_H=\bigcap_{r\ge0}\widehat H^{[r]},
 \qquad \rho_L|_{R_H}:R_H\twoheadrightarrow C.
 \tag{E.43.14}
\]
The upper inclusion of its image in $C$
follows because $S=G_L/C$ is solvable;
the reverse inclusion is the nested-fibre
construction for each $c\in C$.

Set $f=f_0k$.  It remains in
$\widehat F'$ and retains the entire
old source family, including the
identity maps on the full pro-$p$
boundary factors away from two and
three.  In the new source,
$\rho_L(f)=\rho_L(f_\sigma)$.
All $36$ induced coordinates agree,
so the complete new hexagons and
pentagon are those of an actual
Galois element, with the exact based
braid convention of E.38.
The same single $f$ supplies these
maps at every old finite level.

\paragraph{The actual common source and its generation.}
Let $G_{\rm old}$ be an old finite
source containing $P_{20}$.
It is solvable: its factors have
finite solvable geometric kernels
and the solvable quotient $B/H$,
and $\mathscr D$ is solvable as
proved in E.24.
Every common quotient of
$G_{\rm old}$ and $G_L$ therefore
factors through $S$.
The specified map through $P_{20}$
makes $S$ itself their common quotient.
Goursat's lemma identifies their
actual joint $B$ image with
\textup{(E.43.3)}.

The corrected source map is onto
$G_{\rm old}$ by E.42 and onto
$G_L$ because its generator images
are the Galois images followed by
$\rho_L$.  Its two maps to $S$
agree: they evaluate the same
literal source formulas using $f$
and the multiplier one.
The maximal-common-quotient argument
therefore identifies its image with
the full fibre product
\textup{(E.43.3)}.
Projecting this assertion from the
cofinal subsystem proves generation
at every indicated old finite level.
The four-strand comparison is exactly
the one in E.38, since the complete
$D_{23}$ image is preserved and the
new six evaluations are Galois values.
It gives surjectivity onto
$T_{23}\times_{S_L}T_L$, together
with its braid and full-twist identities.

The new program
\texttt{de\_sl2\_boundary\_quotient.py}
retains the exact finite model, all
ten source-generator images, all
$7200$ and $800$ group elements,
and both Sylow projections.
Its comparison uses the proved
E.38 source kernels and an independent
literal order check in the original
$A_5$ permutation representation.
The profinite correction is supplied
by the explicit nested compact fibres
in \textup{(E.43.12)--(E.43.13)}.
No unevaluated Galois word is presented
as a computed literal certificate.
The conclusion concerns this specified
family of finite sources; cofinality
among all finite quotients of $B$
has not been proved.
\end{proof}
